\begin{algorithm}[H]
\SetAlgoLined
\setstretch{0.92}
\SetAlFnt{\normalsize}
\SetArgSty{textnormal}
\SetAlgoNlRelativeSize{-1}
\SetInd{0.55em}{1.05em}
\SetKwInput{KwIn}{Input}
\SetKwInput{KwInit}{Initialize}
\SetKwInput{KwOut}{Output}
\SetKwFor{For}{for}{do}{end for}
\SetKwIF{If}{ElseIf}{Else}{if}{then}{else if}{else}{end if}
\SetKw{KwRet}{return}
\caption{问题三四参数Airy反演与SiC多光束影响检验算法}
\label{alg:q3-airy}
\KwIn{附件3、4硅光谱，问题二SiC厚度及对应光学参数}
\KwInit{固定双振子背景参数与有效质量，设置四参数边界$\Omega_3$和$6$组确定性初值}
\For{$k\in\{10^\circ,15^\circ\}$}{
将反射率百分数转为无量纲小数，并截取$1500$--$3500\,\mathrm{cm}^{-1}$\;
建立候选集合$\mathcal C_k\leftarrow\varnothing$\;
\For{每一组初值$\boldsymbol s\in\mathcal S_3$}{
由双振子--Drude模型计算复折射率$\widetilde n_2$\;
构造Airy多光束反射率$R_{\mathrm{Airy}}$\;
对$(d,n_3,q_N,q_\Gamma)$求解有界非线性最小二乘\;
\If{被动传播条件$0<|P|\le1$不满足}{
舍弃该候选\tcp*[r]{物理可行性检查}
}
记录候选厚度、RMSE、$R^2$及边界活跃参数并写入$\mathcal C_k$\;
}
取$\boldsymbol p_k^*\leftarrow\arg\min_{\boldsymbol p\in\mathcal C_k}\mathrm{RMSE}(\boldsymbol p)$\;
记录$d_k\leftarrow d(\boldsymbol p_k^*)$\;
在$\boldsymbol p_k^*$处构造残差Jacobian并进行尺度化SVD\;
在同一$\boldsymbol p_k^*$下计算双光束截断反射率\;
以$\Delta\mathbf R=\mathbf R_{\mathrm{Airy}}-\mathbf R_{\mathrm{double}}$计算一阶厚度偏移\;
}
计算$d_{\mathrm{Si}}\leftarrow(d_{10^\circ}+d_{15^\circ})/2$\;
将SiC厚度及光学参数代入$\eta_3$，计算两个角度的高阶单束反射量级\;
结合SiC拟合结果与$\eta_3$数量级给出$d_{\mathrm{SiC}}$\;
\KwRet $d_{\mathrm{Si}}$、$d_{\mathrm{SiC}}$及Airy拟合、SVD、高阶反射与模型误差传播结果\;
\KwOut{硅外延层厚度、SiC厚度及多光束影响分析结果}
\end{algorithm}